\documentclass[12pt]{article}
\usepackage[utf8]{inputenc}
\usepackage[T1]{fontenc}
\usepackage{hyperref}
\title{Gundungurra Language Resources for K-12 Curriculum Development}
\author{Compiled by Claude Sonnet 5.5 (claude-sonnet-5-5) for Andr\'as Kornai}
\date{}
\begin{document}
\maketitle

\section{Introduction}
Gundungurra (ISO 639-3 xrd) is an Australian Aboriginal language of the Blue Mountains and Southern Highlands region of New South Wales. It is closely related to Ngunnawal. Sources disagree on its status: Wikipedia mentions revival efforts and an estimated figure for speakers, while Glottolog lists it as extinct. This report rests on a single search round, and no school programme or dedicated Gundungurra revival project turned up.

\section{Primary Educational Resources}
None were found. The Ngunnawal community's revival work on the closely related language may be a useful model, but the search did not show it to be Gundungurra material.

\section{Linguistic and Reference Materials}
\begin{itemize}
\item Wikipedia, Ngunnawal--Gundungurra language: \url{https://en.wikipedia.org/wiki/Ngunnawal%E2%80%93Gundungurra_language}
\item R. H. Mathews, \textit{The Gundungurra Language} (1901), Project Gutenberg: \url{https://gutenberg.org/cache/epub/21866/pg21866-images.html}
\item Same work on ManyBooks: \url{https://manybooks.net/titles/mathewsr2186621866.html}
\item Glottolog record for Kohen 1993, a dictionary of the Gundungurra language (bibliographic record only): \url{https://glottolog.org/resource/reference/id/303015}
\item Glottolog language record: \url{https://glottolog.org/resource/languoid/iso/xrd}
\end{itemize}

\section{Religious and Community Resources}
None were found.

\section{Digital and Audio Resources}
No audio or digital teaching resources were found. A word list page appeared in the results (\url{https://new.marapco.com/dkb39/gundungurra-word-list}), but its reliability is unverified.

\section{Curriculum Development Context}
Documentation is historical: Mathews's 1901 grammar sketch and Kohen's 1993 dictionary. Any revival curriculum would need to work from these sources with the descendant community, and its content would be limited. It should be developed with local Gundungurra and Ngunnawal people, and with a language centre if one exists. The search did not identify one.

\end{document}